\section{Aprobación y Perfiles de tu Personal Interno}

Además de la cuenta individual genérica, tus empleados necesitan tener un perfil de puesto para que funcione su sistema.

\subsection{Doctores Radiólogos y su Firma Digital}

Cuando contrates a un nuevo doctor radiólogo para interpretar estudios, lo primero que harás es crearle su "Usuario" normal con un correo y contraseña (explicado en el capítulo anterior). Después:

\begin{enumerate}
    \item En la pantalla principal, busca la sección \textbf{DoctorsDashboard} y haz clic en \textbf{Reporting doctors}.
    \item Clic en el botón gris de arriba a la derecha que dice \textbf{"Add Reporting doctor"}.
    \item Te pedirá que vincules este puesto con un Usuario. En la lista, busca el correo del doctor recién creado.
    \item Escribe todos sus datos demográficos (Teléfono, Especialidad, Universidad, etc.) requeridos por ley.
\end{enumerate}

\textbf{Cargar la Firma Digital:}
Hay una casilla específica del perfil llamada \textit{Signature} (Firma). Como medida estricta de seguridad contra fraudes e impersonaciones, el radiólogo \textbf{nunca podrá subir esto por sí mismo}. Deberás escanear o conseguir digitalmente su firma en formato PNG, y cargar el archivo aquí. A partir de ese momento, la firma sellará automáticamente todos los PDF que él dicte.

\begin{figure}[H]
    \centering
    \includegraphics[width=0.45\linewidth]{screenshots/admin_signature_upload.png}
    \caption{El administrador adjuntando la fotografía sellada del médico.}
\end{figure}

\subsection{Recepcionistas y Cajeros (Asistentes)}

Al igual que el proceso anterior, para crear asistentes:
\begin{enumerate}
    \item Busca la categoría \textbf{AssistantDashboard} en tu menú inicial.
    \item Ingresa a \textbf{Assistants} y clica en Añadir (\textit{Add}).
    \item Empareja este nombramiento con el correo electrónico del empleado de recepción u oficina general.
    \item Rellena sus datos.
\end{enumerate}

Este apartado es útil para realizar auditorías futuras y saber con certeza qué asistente exacto en turno fue la persona responsable de admitir a pacientes específicos o cobrar ciertos tickets.
